\documentclass[tikz,border=6pt]{standalone}
\usepackage{mathdiagrams}
% algebra_shaded_region — Mock P2 Q6a
% outer (5x+8) × (6x+2), inner unshaded (x+6) × 5x
% Contract: no white boxes; length labels sit along the side they measure.
% Check: (5x+8)(6x+2) − (x+6)(5x) = 30x²+58x+16 − 5x²−30x = 25x²+28x+16.
\begin{document}
\begin{tikzpicture}[md]
  % outer shaded rectangle
  \draw[md accent fill] (0,0) rectangle (5.8,3.6);
  % inner unshaded rectangle
  \draw[md shape] (1.3,0.8) rectangle (4.3,2.7);

  % outer: width above, height at the right
  \node[font=\sffamily\small, text=mdInk] at (2.9,3.88) {$(5x+8)$};
  \node[font=\sffamily\small, text=mdInk] at (6.25,1.8) {$(6x+2)$};

  % inner: width below the inner top edge, height on the inner right edge
  \node[font=\sffamily\small, text=mdInk] at (2.8,2.42) {$(x+6)$};
  \node[font=\sffamily\small, text=mdInk] at (4.55,1.75) {$5x$};

  % equation (data)
  \node[font=\sffamily\small, text=mdInk] at (2.9,-0.55)
    {shaded area $= (5x+8)(6x+2) - (x+6)(5x)$};
\end{tikzpicture}
\end{document}
